\chapter{The Gaussian Chain by Belief Propagation}\label{ch:gaussian-bp}

Chapter~\ref{ch:gaussian} obtained the joint score through the marginal
covariance: forming and inverting the dense $L \times L$ matrix naively costs
$O(L^3)$. That cost is a property of the naive route, not of the problem ---
the posterior precision is tridiagonal, so a banded solver already delivers
the same answer in $O(L)$ --- and this chapter obtains the score by the
message-passing form of exactly that structured solve. The construction is
generic --- it is the sum--product algorithm run on the factor
graph of the posterior --- and for a chain of $K$-state variables it costs
$O(LK^2)$ time and $O(LK)$ memory. For the Gaussian AR(1) chain a message
collapses to two numbers, and the answer is exact; run in
mean--variance arithmetic the recursion is the Kalman filter and the
Rauch--Tung--Striebel smoother, update for update
\citep{kalman1960new, rauch1965maximum}. The chapter closes with the
operational cost of locality: what truncating the evidence range does to the
score, measured exactly. (The closed form of the bulk posterior itself, and
the geometric influence-decay rate $q(\alpha,t)$ it implies, are derived in
Chapter~\ref{ch:em-results}, Section~\ref{sec:em-locality-prediction}, where
they are used as a design calculation.)

\section{The posterior and its factor graph}
\label{sec:bp-chain}\label{sec:bp-problem}

The matrix route is global: one $L \times L$ inversion. The score--posterior
identity \eqref{eq:tweedie-joint} says the same object is a posterior
expectation, and that posterior has a structure --- a Markov chain observed
through independent per-frame noise --- built for local computation.

\subsection{The posterior is a hidden Markov model}

Bayes' rule multiplies the chain prior \eqref{eq:chain-factorisation} by the
per-frame likelihoods \eqref{eq:g-channel}. Nothing in the likelihood couples
two frames --- each $x_k$ sees only its own $a_k$ --- so the product keeps the
prior's two-variable structure:
\begin{equation}
  p(a \,|\, x) \;\propto\;
  \underbrace{\psi_{\mathrm{pr}}(a_0)}_{\Nd(a_0;\,0,\,1)}
  \prod_{k=0}^{L-2}
  \underbrace{\psi^{(k)}_{\mathrm{tr}}(a_k, a_{k+1})}
            _{\Nd(a_{k+1};\,\alpha a_k,\,\sigma_\eta^2)}
  \prod_{k=0}^{L-1}
  \underbrace{\psi^{(k)}_{\mathrm{ob}}(a_k)}
            _{\Nd(x_k;\,e^{-t}a_k,\,\Deltat)} .
  \label{eq:bp-posterior}
\end{equation}
This is a hidden Markov model \citep{rabiner1989tutorial}: a latent chain with
nearest-neighbour transitions, each state reported once through an independent
channel. The data $x_k$ are fixed numbers absorbed into the observation
factors, so $\psi^{(k)}_{\mathrm{ob}}$ is a function of the single variable
$a_k$. By \eqref{eq:tweedie-joint}, the $k$-th component of the joint score is
fixed by one number from this posterior --- the smoothing mean
$\E[a_k \mid x]$.

\subsection{The factor graph is a tree}

The factor graph (Section~\ref{sec:factor-graphs}) of
\eqref{eq:bp-posterior} has
\begin{itemize}
\item $L$ \emph{variable} nodes, one per $a_k$;
\item $2L$ \emph{factor} nodes: one prior $\psi_{\mathrm{pr}}$, $L-1$
  transitions $\psi^{(k)}_{\mathrm{tr}}$, $L$ observations
  $\psi^{(k)}_{\mathrm{ob}}$;
\item $3L - 1$ \emph{edges}: each transition factor joins two variables
  ($2(L-1)$ edges), the prior and each observation factor join one
  ($L + 1$ edges).
\end{itemize}
A connected factor graph with one fewer edge than node ($3L-1$ against $3L$) is
a tree. Its shape is a spine of variables $a_k$ linked by transition factors,
with one observation leaf per spine node
(Figure~\ref{fig:posterior-factor-graph}). Conditioning on the data did not
change the graph of the prior: per-node evidence only re-weights an existing
factor and hangs a leaf, and never joins two spine nodes
(Section~\ref{sec:factor-graphs}).

\section{The sum--product sweep and its cost}
\label{sec:bp-sweep}

\subsection{The two sweeps}

Specialising the sum--product rules \eqref{eq:bp-rules} to this tree, the
schedule collapses to one forward and one backward pass
(Theorem~\ref{thm:tree-exact}).

\begin{proposition}[Chain sweeps]\label{prop:bp-convA}
Define messages
\begin{align}
  \mu_{\to 0}(a_0) &:= \psi_{\mathrm{pr}}(a_0),
  &
  \mu_{\to k+1}(a_{k+1}) &:= \int
    \psi^{(k)}_{\mathrm{tr}}(a_k, a_{k+1})\,
    \psi^{(k)}_{\mathrm{ob}}(a_k)\,
    \mu_{\to k}(a_k)\; \dd a_k ,
  \label{eq:bp-fwd}\\
  \mu_{\leftarrow L-1}(a_{L-1}) &:= 1,
  &
  \mu_{\leftarrow k-1}(a_{k-1}) &:= \int
    \psi^{(k-1)}_{\mathrm{tr}}(a_{k-1}, a_k)\,
    \psi^{(k)}_{\mathrm{ob}}(a_k)\,
    \mu_{\leftarrow k}(a_k)\; \dd a_k .
  \label{eq:bp-bwd}
\end{align}
Then for every $k$,
\begin{equation}
  p(a_k \,|\, x) \;\propto\;
  \mu_{\to k}(a_k)\;
  \psi^{(k)}_{\mathrm{ob}}(a_k)\;
  \mu_{\leftarrow k}(a_k) :
  \label{eq:bp-combine}
\end{equation}
the forward message into $a_k$ carries the evidence
$x_0, \dots, x_{k-1}$ only, the backward message carries
$x_{k+1}, \dots, x_{L-1}$ only, and the local evidence $x_k$ enters once,
at combination time.
\end{proposition}


\begin{toolbox}{Why the two sweeps give the exact smoothing marginal}

Fix a site $k$. Starting from the posterior factorisation
\eqref{eq:bp-posterior}, separate all factors into those lying to the left of
$a_k$, the local observation at $a_k$, and those lying to its right:
\begin{align}
  p(a\mid x)
  \propto\;&
  \underbrace{
  \psi_{\mathrm{pr}}(a_0)
  \prod_{j=0}^{k-1}
    \psi_{\mathrm{tr}}^{(j)}(a_j,a_{j+1})
  \prod_{j=0}^{k-1}
    \psi_{\mathrm{ob}}^{(j)}(a_j)
  }_{\text{left of }a_k}
  \nonumber\\
  &\times
  \psi_{\mathrm{ob}}^{(k)}(a_k)
  \nonumber\\
  &\times
  \underbrace{
  \prod_{j=k}^{L-2}
    \psi_{\mathrm{tr}}^{(j)}(a_j,a_{j+1})
  \prod_{j=k+1}^{L-1}
    \psi_{\mathrm{ob}}^{(j)}(a_j)
  }_{\text{right of }a_k}.
  \label{eq:bp-split-at-k}
\end{align}
To obtain $p(a_k\mid x)$, integrate over every variable except $a_k$.
Because the factor graph is a chain, the variables on the two sides of
$a_k$ occur in disjoint factors once $a_k$ is held fixed. The integral
therefore separates:
\begin{align}
  p(a_k\mid x)
  \propto\;&
  \left[
    \int
    \psi_{\mathrm{pr}}(a_0)
    \prod_{j=0}^{k-1}
      \psi_{\mathrm{tr}}^{(j)}(a_j,a_{j+1})
    \prod_{j=0}^{k-1}
      \psi_{\mathrm{ob}}^{(j)}(a_j)
    \,\dd a_0\cdots\dd a_{k-1}
  \right]
  \nonumber\\
  &\times
  \psi_{\mathrm{ob}}^{(k)}(a_k)
  \nonumber\\
  &\times
  \left[
    \int
    \prod_{j=k}^{L-2}
      \psi_{\mathrm{tr}}^{(j)}(a_j,a_{j+1})
    \prod_{j=k+1}^{L-1}
      \psi_{\mathrm{ob}}^{(j)}(a_j)
    \,\dd a_{k+1}\cdots\dd a_{L-1}
  \right].
  \label{eq:bp-split-integrals}
\end{align}

The first bracket is precisely the information arriving at $a_k$ from the
left. Define it as
\begin{equation}
  \mu_{\to k}(a_k)
  :=
  \int
  \psi_{\mathrm{pr}}(a_0)
  \prod_{j=0}^{k-1}
    \psi_{\mathrm{tr}}^{(j)}(a_j,a_{j+1})
  \prod_{j=0}^{k-1}
    \psi_{\mathrm{ob}}^{(j)}(a_j)
  \,\dd a_0\cdots\dd a_{k-1}.
  \label{eq:bp-fwd-partial}
\end{equation}
For $k=0$ there is nothing to integrate, so
\begin{equation}
  \mu_{\to0}(a_0)=\psi_{\mathrm{pr}}(a_0).
\end{equation}
To obtain the message one step farther to the right, separate the last
integration variable $a_k$:
\begin{align}
  \mu_{\to k+1}(a_{k+1})
  &=
  \int
  \psi_{\mathrm{tr}}^{(k)}(a_k,a_{k+1})\,
  \psi_{\mathrm{ob}}^{(k)}(a_k)\,
  \mu_{\to k}(a_k)\,
  \dd a_k,
\end{align}
which is exactly the forward recursion \eqref{eq:bp-fwd}.

Similarly, define the information arriving from the right by
\begin{equation}
  \mu_{\leftarrow k}(a_k)
  :=
  \int
  \prod_{j=k}^{L-2}
    \psi_{\mathrm{tr}}^{(j)}(a_j,a_{j+1})
  \prod_{j=k+1}^{L-1}
    \psi_{\mathrm{ob}}^{(j)}(a_j)
  \,\dd a_{k+1}\cdots\dd a_{L-1}.
  \label{eq:bp-bwd-partial}
\end{equation}
At the right endpoint there are no factors beyond $a_{L-1}$, so
\begin{equation}
  \mu_{\leftarrow L-1}(a_{L-1})=1.
\end{equation}
Peeling off the first variable on the right gives
\begin{align}
  \mu_{\leftarrow k-1}(a_{k-1})
  &=
  \int
  \psi_{\mathrm{tr}}^{(k-1)}(a_{k-1},a_k)\,
  \psi_{\mathrm{ob}}^{(k)}(a_k)\,
  \mu_{\leftarrow k}(a_k)\,
  \dd a_k,
\end{align}
which is the backward recursion \eqref{eq:bp-bwd}.

Substituting the two partial marginalisations into
\eqref{eq:bp-split-integrals} yields
\begin{equation}
  \boxed{
  p(a_k\mid x)
  \propto
  \mu_{\to k}(a_k)\,
  \psi_{\mathrm{ob}}^{(k)}(a_k)\,
  \mu_{\leftarrow k}(a_k)
  }.
  \label{eq:bp-combine-proof}
\end{equation}
The separation of the evidence is now explicit:
\begin{align}
  \mu_{\to k}
  &\;\text{depends only on}\;
  x_0,\ldots,x_{k-1},\\
  \psi_{\mathrm{ob}}^{(k)}
  &\;\text{contains only}\;
  x_k,\\
  \mu_{\leftarrow k}
  &\;\text{depends only on}\;
  x_{k+1},\ldots,x_{L-1}.
\end{align}
Thus every observation enters the marginal exactly once. Because the factor
graph is a tree, no approximation has been made.
\hfill$\square$
\end{toolbox}

\paragraph{Computational cost.}
For a discrete chain in which each $a_k$ takes $K$ values, every message is a
vector in $\R^K$. Consider, for example, the forward update
\[
  \mu_{\to k+1}(a_{k+1})
  =
  \sum_{a_k}
  \psi_{\mathrm{tr}}^{(k)}(a_k,a_{k+1})\,
  \psi_{\mathrm{ob}}^{(k)}(a_k)\,
  \mu_{\to k}(a_k).
\]
For each of the $K$ possible values of $a_{k+1}$, the sum runs over the $K$
possible values of $a_k$, so one transition message costs $O(K^2)$. A
forward--backward sweep sends $O(L)$ such messages along the chain; therefore
\begin{equation}
  \boxed{
  \text{time}=O(LK^2).
  }
  \label{eq:bp-time}
\end{equation}
Each stored message contains $K$ numbers. Keeping the messages required to
construct all $L$ smoothing marginals therefore costs
\begin{equation}
  \boxed{
  \text{space}=O(LK).
  }
  \label{eq:bp-space}
\end{equation}
Thus sum--product is linear in the sequence length $L$; the $K^2$ factor is
the local cost of marginalising across one pairwise transition factor.



\section{The Gaussian closure}
\label{sec:bp-gaussian}

Two facts close the recursion: a product of
Gaussians in the same variable is Gaussian, and an
affine map of a Gaussian, marginalised, is Gaussian.
The observation factor, read as a function of the clean variable, is itself
Gaussian,
\begin{equation*}
  \psi^{(k)}_{\mathrm{ob}}(a_k)
  \;\propto\; \exp\!\Big(-\tfrac{e^{-2t}}{2\Deltat}\,(a_k - e^{t} x_k)^2\Big)
  \;=\; \Nd\big(a_k;\, e^{t} x_k,\, \Deltat\, e^{2t}\big),
\end{equation*}
a pseudo-observation at the deconvolved location $e^t x_k$ with precision
$e^{-2t}/\Deltat$.


\begin{theorem}[Gaussian closure of BP on the chain]
\label{thm:bp-closure}
For the Gaussian AR(1) chain observed through the OU channel, every
sum--product message and every one-node belief is Gaussian. A message can
therefore be represented by two scalar natural parameters, and each message
update requires $O(1)$ arithmetic. Since a forward--backward sweep contains
$O(L)$ messages, all posterior marginals, and hence the complete joint score,
are obtained exactly in
\begin{equation}
  \boxed{
  \text{time}=O(L),
  \qquad
  \text{space}=O(L)
  }
\end{equation}
when all smoothing marginals are retained.
\end{theorem}

\begin{toolbox}{Why Gaussian BP closes exactly}

It is convenient to write an unnormalised scalar Gaussian in natural form,
\begin{equation}
  \mu(a)
  \propto
  \exp\!\left(
    -\frac12 \lambda a^2+r a
  \right),
  \label{eq:bp-natural-message}
\end{equation}
where $\lambda$ is the precision and $r$ the field. When $\lambda>0$,
\begin{equation}
  v=\frac{1}{\lambda},
  \qquad
  m=\frac{r}{\lambda}.
\end{equation}
Thus a complete message is represented by only the pair $(\lambda,r)$.

The observation factor at site $k$ is itself Gaussian as a function of
$a_k$:
\begin{align}
  \psi_{\mathrm{ob}}^{(k)}(a_k)
  &=
  \Nd(x_k;e^{-t}a_k,\Deltat)\\
  &\propto
  \exp\!\left[
    -\frac{(x_k-e^{-t}a_k)^2}{2\Deltat}
  \right]\\
  &\propto
  \exp\!\left[
    -\frac12
      \frac{e^{-2t}}{\Deltat}a_k^2
    +
      \frac{e^{-t}x_k}{\Deltat}a_k
  \right].
  \label{eq:bp-observation-natural}
\end{align}
Hence its natural parameters are
\begin{equation}
  \lambda_{\mathrm{ob}}
  =
  \frac{e^{-2t}}{\Deltat},
  \qquad
  r_{\mathrm{ob},k}
  =
  \frac{e^{-t}x_k}{\Deltat}.
  \label{eq:bp-observation-parameters}
\end{equation}

\medskip
\noindent\textbf{1. Products preserve Gaussianity.}
Suppose an incoming message has parameters $(\lambda,r)$. Multiplying it by
the local observation factor gives
\begin{align}
  \mu(a_k)\psi_{\mathrm{ob}}^{(k)}(a_k)
  &\propto
  \exp\!\left[
    -\frac12
      (\lambda+\lambda_{\mathrm{ob}})a_k^2
    +
      (r+r_{\mathrm{ob},k})a_k
  \right].
\end{align}
Thus the product remains Gaussian, with
\begin{equation}
  \bar\lambda
  =
  \lambda+\lambda_{\mathrm{ob}},
  \qquad
  \bar r
  =
  r+r_{\mathrm{ob},k}.
  \label{eq:bp-product-update}
\end{equation}
In natural parameters, multiplying Gaussian factors is therefore just
addition.

\medskip
\noindent\textbf{2. Passing through the transition preserves Gaussianity.}
The AR(1) transition is
\begin{equation}
  a_{k+1}
  =
  \alpha a_k+\eta_k,
  \qquad
  \eta_k\sim\Nd(0,\sigma_\eta^2).
\end{equation}
After incorporating the observation at $k$, the variable $a_k$ has Gaussian
parameters $(\bar\lambda,\bar r)$, hence
\begin{equation}
  a_k
  \sim
  \Nd\!\left(
    \frac{\bar r}{\bar\lambda},
    \frac{1}{\bar\lambda}
  \right).
\end{equation}
An affine transformation of a Gaussian plus independent Gaussian noise is
again Gaussian, so
\begin{align}
  a_{k+1}
  &\sim
  \Nd\!\left(
    \alpha\frac{\bar r}{\bar\lambda},
    \sigma_\eta^2+\frac{\alpha^2}{\bar\lambda}
  \right).
\end{align}
Therefore the outgoing forward message is Gaussian with natural parameters
\begin{align}
  \lambda_{\to k+1}
  &=
  \frac{\bar\lambda}
       {\alpha^2+\sigma_\eta^2\bar\lambda},
  \label{eq:bp-forward-precision}\\
  r_{\to k+1}
  &=
  \frac{\alpha\bar r}
       {\alpha^2+\sigma_\eta^2\bar\lambda}.
  \label{eq:bp-forward-field}
\end{align}
The map
\[
  (\lambda_{\to k},r_{\to k})
  \longmapsto
  (\lambda_{\to k+1},r_{\to k+1})
\]
therefore uses only a fixed number of scalar arithmetic operations.

The backward recursion is closed for exactly the same reason. If the
observation at $k$ and the incoming right message combine to
$(\bar\lambda,\bar r)$, then integrating
\begin{equation}
  \int
  \Nd(a_k;\alpha a_{k-1},\sigma_\eta^2)
  \exp\!\left(
    -\frac12\bar\lambda a_k^2+\bar r a_k
  \right)
  \dd a_k
\end{equation}
gives another Gaussian function of $a_{k-1}$, with
\begin{align}
  \lambda_{\leftarrow k-1}
  &=
  \frac{\alpha^2\bar\lambda}
       {1+\sigma_\eta^2\bar\lambda},
  \label{eq:bp-backward-precision}\\
  r_{\leftarrow k-1}
  &=
  \frac{\alpha\bar r}
       {1+\sigma_\eta^2\bar\lambda}.
  \label{eq:bp-backward-field}
\end{align}

\medskip
\noindent\textbf{3. Combining the two sweeps preserves Gaussianity.}
At node $k$ the exact smoothing belief is
\[
  p(a_k\mid x)
  \propto
  \mu_{\to k}(a_k)\,
  \psi_{\mathrm{ob}}^{(k)}(a_k)\,
  \mu_{\leftarrow k}(a_k).
\]
All three factors are Gaussian, so their natural parameters add:
\begin{align}
  \lambda_k^{\mathrm{post}}
  &=
  \lambda_{\to k}
  +\lambda_{\mathrm{ob}}
  +\lambda_{\leftarrow k},
  \label{eq:bp-belief-precision}\\
  r_k^{\mathrm{post}}
  &=
  r_{\to k}
  +r_{\mathrm{ob},k}
  +r_{\leftarrow k}.
  \label{eq:bp-belief-field}
\end{align}
Consequently
\begin{equation}
  \boxed{
  \E[a_k\mid x]
  =
  \frac{r_k^{\mathrm{post}}}
       {\lambda_k^{\mathrm{post}}}
  }.
  \label{eq:bp-belief-mean}
\end{equation}
Substitution into the score--posterior identity gives
\begin{equation}
  \score_k(x,t)
  =
  \frac{
    e^{-t}r_k^{\mathrm{post}}/
    \lambda_k^{\mathrm{post}}
    -x_k
  }{\Deltat}.
\end{equation}

The closure is exact: no Gaussian approximation has been introduced.
Gaussian factors are closed algebraically under the two operations used by
sum--product, multiplication and marginalisation, and the factor graph is a
tree.
\hfill$\square$
\end{toolbox}

\begin{remark}[Why not TAP or approximate message passing?]
\label{rmk:bp-no-amp}
The Gaussian closure above should not be confused with the approximations
used in TAP or approximate message passing (AMP). TAP equations originate
as mean-field equations for densely interacting systems and introduce an
Onsager reaction correction to account for correlations neglected by naive
mean-field theory. AMP turns closely related ideas into an iterative
message-passing algorithm for high-dimensional, typically dense random
models: when each variable receives many individually weak contributions,
the full collection of BP messages can be compressed into approximately
Gaussian scalar updates together with an Onsager correction.

None of those approximations is needed here. Our factor graph is a
one-dimensional tree of finite degree, so ordinary sum--product is already
exact. Moreover, in the Gaussian model the messages are Gaussian because
Gaussian multiplication and marginalisation are exact algebraic operations,
not because many weak contributions become Gaussian in a large-system
limit. Using TAP or AMP would therefore replace an exact algorithm
with an approximation designed for a different, typically dense and loopy,
regime.
\end{remark}

\paragraph{Numerical agreement of the two routes.} The audit suite runs BP
against the matrix route over $80$ configurations
($L \in \{2,3,5,8,16\}$, $\alpha \in \{0.2, 0.5, 0.9, -0.4\}$,
$t \in \{0.05, 0.3, 1, 3\}$, random $x$): maximum disagreement
$1.3 \times 10^{-15}$ on posterior means, $3.9 \times 10^{-15}$ on variances,
$1.2 \times 10^{-14}$ on scores --- double-precision rounding accumulated along
the recursion. An independent implementation, written from scratch against the
recursion specification alone, reproduces the same agreement
(Appendix~\ref{app:repro}). One quantitative detail from a worked $L = 3$
instance: at $x = (1.2, -0.4, 0.8)$ with $\alpha = 0.9$, $t = 0.05$ the
score component $S_1$ is $+5.71$, fourteen times the $-x_1 = +0.40$ that a
per-frame marginal score would give. The neighbours, not the frame itself,
set the size of the correction; it is inter-frame evidence sharing in
numbers.

\section{The price of locality}
\label{sec:bp-locality}\label{sec:g-bulk}\label{sec:g-locality}\label{sec:bp-truncation}

Belief propagation on the chain is exact, but it is not local: the message
reaching frame $k$ has travelled the whole sequence. Chapter~\ref{ch:gaussian}
showed that the exact score operator $\Qt$ is tridiagonal only at $t = 0$ and
dense thereafter, and measured how much coupling mass sits outside the Markov
band (Figure~\ref{fig:g-markov-survival}). That is a statement about a matrix.
This section asks the operational question behind it: \emph{what does it cost
to ignore that mass?} --- how much of the score survives if evidence is only
allowed to travel a few frames.

\subsection{Two ways to go local}

\paragraph{Truncate the answer.} Take the exact operator $\Qt$, keep only the
entries within a band, and use what is left:
\begin{equation}
  \widehat S^{(b)}(x,t) \;=\; -\,\Qt^{(b)} x,
  \qquad
  \big(\Qt^{(b)}\big)_{ij} = (\Qt)_{ij}\,\mathbf{1}\{|i-j| \le b\}.
  \label{eq:bp-banded}
\end{equation}
At $b = 1$ this is the tridiagonal --- Markov --- truncation of the score.

\paragraph{Truncate the algorithm.} Let the messages of
Proposition~\ref{prop:bp-convA} travel at most $r$ hops before frame $k$ is
read off: beyond the cut, the incoming message is replaced by what it was
before any evidence arrived, the stationary prior marginal. Because a
contiguous window of a stationary AR(1) chain is itself a stationary AR(1)
chain, this is not an approximation of the window's inference but exact
inference on the window,
\begin{equation}
  m^{(r)}_k \;=\; \E\big[a_k \,\big|\, x_{k-r}, \dots, x_{k+r}\big],
  \qquad
  \widehat S^{(r)}_k(x,t) \;=\; \frac{e^{-t} m^{(r)}_k - x_k}{\Deltat},
  \label{eq:bp-windowed}
\end{equation}

Both approximations are linear in the noisy trajectory. Writing either of
them as
\begin{equation}
  \widehat S(x,t)=-Mx,
\end{equation}
while the exact Gaussian score is
\begin{equation}
  S(x,t)=-\Qt x,
\end{equation}
their error on a particular trajectory is
\begin{equation}
  \widehat S(x,t)-S(x,t)
  =
  -(M-\Qt)x.
\end{equation}
Its size depends on the particular realization of $x$, so we average the
squared error over the actual noisy data law
$x\sim\Nd(0,\Sigt)$. To make errors comparable across diffusion times and
locality budgets, we normalize by the average squared magnitude of the exact
score itself. This gives the relative RMS score error
\begin{equation}
  \varepsilon(M)
  :=
  \sqrt{
    \frac{
      \E\bigl\|(M-\Qt)x\bigr\|^2
    }{
      \E\bigl\|\Qt x\bigr\|^2
    }
  }.
  \label{eq:bp-relerr-def}
\end{equation}
Thus $\varepsilon(M)=0$ means an exact score, while, for example,
$\varepsilon(M)=0.1$ means that the RMS approximation error is $10\%$ of the
RMS size of the exact score.

Because both estimators are linear, the expectations can be evaluated
without Monte Carlo sampling. For any matrix $A$ and any centred random
vector $x$ with covariance $\Sigt$,
\begin{align}
  \E[x^{\T}Ax]
  &=
  \E\!\left[\operatorname{tr}(Axx^{\T})\right] \\
  &=
  \operatorname{tr}\!\left(A\,\E[xx^{\T}]\right) \\
  &=
  \operatorname{tr}(A\Sigt).
  \label{eq:bp-quadratic-trace}
\end{align}
Applying this identity with
$A=(M-\Qt)^{\T}(M-\Qt)$ gives
\begin{align}
  \E\bigl\|(M-\Qt)x\bigr\|^2
  &=
  \operatorname{tr}\!\left[
    (M-\Qt)^{\T}(M-\Qt)\Sigt
  \right] \\
  &=
  \operatorname{tr}\!\left[
    (M-\Qt)\Sigt(M-\Qt)^{\T}
  \right],
  \label{eq:bp-error-numerator}
\end{align}
where the second equality uses cyclicity of the trace.

The denominator simplifies further because
$\Qt=\Sigt^{-1}$ and both matrices are symmetric:
\begin{align}
  \E\|\Qt x\|^2
  &=
  \operatorname{tr}(\Qt\Sigt\Qt) \\
  &=
  \operatorname{tr}(\Qt).
  \label{eq:bp-error-denominator}
\end{align}
Hence
\begin{equation}
  \boxed{
  \varepsilon(M)
  =
  \sqrt{
    \frac{
      \operatorname{tr}\!\left[
        (M-\Qt)\Sigt(M-\Qt)^{\T}
      \right]
    }{
      \operatorname{tr}\Qt
    }
  }.
  }
  \label{eq:bp-relerr}
\end{equation}

Equation~\eqref{eq:bp-relerr} is therefore a population error, not a
sampled estimate: once $M$, $\Sigt$, and $\Qt$ are known, the expected score
error is deterministic.


\subsection{What the neighbours are worth}

Table~\ref{tab:g-truncation} evaluates \eqref{eq:bp-relerr} for both families.

\begin{center}
\begin{minipage}{\linewidth}\centering
\captionof{table}[The price of locality]{Relative RMS score error
\eqref{eq:bp-relerr} of the two truncations ($L = 40$, $\alpha = 0.8$). Both
estimators are linear in $x$, so every entry is closed-form, not sampled. The
last two rows are derived: the fraction of the no-neighbour error removed by
admitting a single hop, and the penalty for truncating the answer rather than
the algorithm at equal budget.}
\label{tab:g-truncation}
\begin{tabular}{@{}lcccc@{}}
\toprule
relative RMS score error & $t = 0.05$ & $t = 0.3$ & $t = 1$ & $t = 3$ \\
\midrule
no neighbours, $r = 0$        & $0.815$ & $0.567$ & $0.194$ & $0.0045$ \\
one hop, $r = 1$              & $0.123$ & $0.235$ & $0.130$ & $0.0035$ \\
two hops, $r = 2$             & $0.020$ & $0.107$ & $0.089$ & $0.0028$ \\
\addlinespace
banded answer, $b = 1$        & $0.210$ & $0.301$ & $0.135$ & $0.0035$ \\
banded answer, $b = 2$        & $0.032$ & $0.134$ & $0.092$ & $0.0028$ \\
\midrule
removed by the first hop      & $85\%$  & $58\%$  & $33\%$  & $21\%$ \\
penalty, $b{=}1$ over $r{=}1$ & $1.71\times$ & $1.28\times$ & $1.03\times$ & $1.00\times$ \\
\bottomrule
\end{tabular}
\end{minipage}
\end{center}

The first row is the cost of treating the frames as independent. At
$t = 0.05$ the isolated-frame score is wrong by $81.5\%$ in relative RMS: almost
all of the score at small diffusion times \emph{is} inter-frame structure,
which is the quantitative form of the point made with the marginal score in
Section~\ref{sec:g-score}. This error does not vanish as $t \to 0$: the
$r = 0$ estimator converges to $-x$ while the exact score converges to
$-\Qz x$, so its low-noise limit is set by the clean coupling itself. The second row is what one hop of message passing
recovers: the error falls to $0.123$, so a single nearest-neighbour exchange
--- the Markov band, and nothing more --- removes $85\%$ of it. Two hops
remove $98\%$.

That fraction falls with diffusion time ($85\%, 58\%, 33\%, 21\%$), which
sounds like the Markov band mattering less; the first row shows why it does
not. By $t = 3$ the entire inter-frame content of the score is worth $0.0045$
in relative error, so a model that ignores every neighbour is already accurate
to half a percent and there is nothing left for any hop to recover. The band
carries the largest share of the coupling exactly where the coupling is large.

\begin{figure}[t]\centering
\includegraphics[width=\linewidth]{fig_truncation.pdf}
\caption[The price of locality]{The two truncations ($L = 40$,
$\alpha = 0.8$; every value closed-form by \eqref{eq:bp-relerr}, none
sampled). \emph{Left:} relative RMS score error against diffusion time, for
message range $r$ (solid) and matrix band $b$ (dashed). For $r \geq 1$ both
are worst at intermediate $t$ and vanish at both ends --- at $t \to 0$
because the clean score is tridiagonal, so any window that retains the Markov
blanket becomes exact; the no-neighbour estimator $r = 0$ is the exception,
since the clean score of frame $k$ genuinely uses its neighbours whenever
$\alpha \neq 0$, and its error tends to a positive limit instead. The gap
between the families closes as $t$ grows. \emph{Right:} the same error against the locality budget
at three diffusion times: geometric decay, at a rate that worsens with $t$,
with the banded answer (squares) tracking above the windowed inference
(circles) at every budget.}
\label{fig:truncation}
\end{figure}

